\scanpage{59}
\mspage{65}

% slanted
\margin{Meilleure terminologie : \unsure{$f$ est|fonct} $\unsure{\mathbb D|\underline D}$-\uline{asphérique} ($\mathbb D$-équiv stricte), ou $\mathbb D$-\uline{coasphérique} ($\mathbb D$-équiv. co-stricte)}

\uline{Définition 3} \inserted{Soit $\mathbb D$ un prédérivateur (provisoir$^{\mathrm{mt}}$ un dér. à g. ou à dr.)} On dit que $f$ est une $\mathbb D$-équiv. \uline{stricte}, \deleted{\unsure{ou}} \deleted{que $f$ est un foncteur $\mathbb D$-cofinal}, si $f$ satisfait les conditions \inserted{a)} \deleted{équivalentes} de prop. 2. On dit que c'est une $\mathbb D$-équiv costricte si $f$ satisfait celles de prop. $2'$ ; on dit aussi parfois que $f$ est $\mathbb D$-cofinal. (la condition duale, i.e. $\mathbb D$-équiv. stricte, s'appellerait co-$\mathbb D$-cofinal !)

\uline{Prop. 3} \deleted{\unsure{Soit}} \unsure{Soit} $X\longrightarrow Y\longrightarrow Z$. Si $f,g$ sont des $\mathbb D$-équivalences strictes, alors $gf$ aussi. Si $gf$ et $f$ sont des $\mathbb D$-équiv. strictes, alors $g$ aussi.

\inserted{si $\mathbb D$ dér. à g. ou à dr.} \deleted{Toute \deleted{\unsure{mor}} \inserted{équivalence} est une $\mathbb D$-équivalence stricte (en effet $f:X\longrightarrow Y$ une équiv., alors $\forall y\in Y$, $X_{/y}$ est équiv. à la cat. ponctuelle, et on utilise que une \uline{cat. équiv. à la cat. ponctuelle est $\mathbb D$-contractile})}

% slanted
\margin{\unsure{Cas où} $f$ \unsure{a} un adj. à droite.}

% slanted
\margin{$\swarrow$ cf cor. 1 de Dér. 5 (g)}

On utilise le fait que la notion d'isom. local est stable par chgt. de base.

\deleted{\uline{Prop. 4} \inserted{($\mathbb D$ dér. à g.)} Soit $f:X\longrightarrow Y$ une $\mathbb D\unsure{^{+}|^{*}}$-équiv. costricte.}

\diagram{\begin{tikzcd}[ampersand replacement=\&,row sep=tiny]
\deleted{\text{Alors }f^*:\mathbb D(Y)} \arrow[rr,dash] \arrow[dr,dotted,"p"'] \& {} \& \deleted{\mathbb D(X)} \arrow[dl,dotted,"q"] \arrow[r,phantom] \& \deleted{\text{est pl. fidèle}} \\
{} \& \deleted{e} \& {} \& {}
\end{tikzcd}}

\deleted{Alors pour $F\in\mathrm{Ob}\,\mathbb D(Y)$, \inserted{(NB $F$ \uline{\unsure{non|pas} néc. constant})} le morphisme \deleted{\unsure{de}}}

\[ \deleted{q_*(F)\xrightarrow{\sim}p_*(f^*F)\quad\text{i.e. }H^\bullet_{\mathbb D}(Y,F)\xrightarrow{\sim}H^\bullet_{\mathbb D}(X,f^*F)} \]

\deleted{est un iso.}

\deleted{C'est un ``rappel'' de l'axiome Dér 5 (g) pour les dérivateurs à gauche. Si $Z$ est catégorie avec objet final $e$, alors appliquant \uline{Der} 5 \unsure{g)}, à $\{e\}\longrightarrow Z$, on trouve que pour $F$ sur $Z$, on a $H^\bullet_{\mathbb D}(Z,F)\simeq F_e$, ce qui (\unsure{prenant|posant} $F$ constant) implique que $Z$ est $\mathbb D$-contractile.}
% archive sheet number: 58
